\scanpage{5}
% L1
a fortiori un topos.

% L2
\deleted{\illegible{1}}

% L3
On va définir, pour deux morphismes de topos

% L4
\[
\deleted{\unsure{f}}u\colon E\longrightarrow E',\qquad \deleted{\unsure{g}}v\colon F\to F'
\]

% L5
un morphisme de topos

% L6
\[
\Pi(u,v)\colon\pi(E,F)\longrightarrow\pi(E',F')
\]

% L7–8
dont l'interprétation intrinsèque symétrique est ainsi, en termes des ``bifaisceaux'' $E^o\times F^o\longrightarrow(\mathrm{Ens})$, et $E'{}^o\times F'{}^o\longrightarrow(\mathrm{Ens})$,

% L9–11
\[
\begin{aligned}
&\deleted{\unsure{\pi(u,v)^*(\varphi')(x,y)=\varphi'(u_*(x),v_*(y))}}\\
&\Pi(u,v)_{\ast}^{\deleted{\ast}}(\varphi)(X',Y')=\varphi(u^*X',v^*Y')\quad:\pi(E,F)\longrightarrow\pi(E',F')\\
&\pi(u,v)^*(\varphi)(X,Y)\overset{?}{=}
\end{aligned}
\]

% L12–13
Pour prouver que \deleted{\unsure{c'est associé}} foncteur précédent est associé à un morphisme de topos, nous sommes ramenés, en \deleted{\unsure{fa|le}} factorisant en

% L14
\[
\pi(E,F)\longrightarrow\pi(E,F')\longrightarrow\pi(E',F')
\]

% L15–16
au cas où l'un des deux morphismes $u,v$ est l'identité ; soit $u$, donc à étudier

% L17
\[
\pi(\mathrm{id}_E,v)_*\colon\pi(E,F)\longrightarrow\pi(E,F')
\]

% L18–21
On peut supposer \deleted{$v\colon F\to F'$ défini par un morphisme de \inserted{petits} sites $v_0^*\colon C'\to C$, $C,C'$ stables par \deleted{\unsure{produits}} $\varprojlim$ finies, leur top. moins fine que la top can. \deleted{Donc} On a \deleted{\unsure{évidemment}} alors un diagramme commutatif \unsure{$\pi(\mathrm{id}_E,v)_*$} de topos}

% L22–25-left
\diagram{\begin{tikzcd}[ampersand replacement=\&,column sep=large,execute at end picture={\draw (current bounding box.north west) -- (current bounding box.south);}]
\deleted{F}\arrow[r,"{\deleted{\unsure{v_0|\widetilde v_0}}}"] \& \deleted{F'}\\
\deleted{\widetilde C=F}\arrow[r,"{\deleted{\widetilde v_0}}"]\arrow[d,hook] \& \deleted{\widetilde C'=F'}\arrow[d,hook]\\
\deleted{\widehat C}\arrow[r,"{\deleted{\widehat v_0}}"] \& \deleted{\unsure{\widehat C'|C'}}
\end{tikzcd}}

% L26-left
\deleted{et par suite un diagr\unsure{.}}

% L22–25-right
$E=\widetilde C$, et \deleted{\unsure{on a}} alors $\pi(\mathrm{id}_E,v)_*$ s'insère dans un diagramme \inserted{\unsure{ess.}} commutatif

% L26–29
\diagram{\begin{tikzcd}[ampersand replacement=\&,column sep=huge,row sep=large]
\pi(\widetilde C,F)\arrow[r,"{\pi(\mathrm{id}_{\widetilde C},v)_*}"]\arrow[d,hook,"{\pi(i_C,\mathrm{id}_F)_*}"'] \& \pi(\widetilde C,F')\arrow[d,hook,"{\pi(i_C,\mathrm{id}_{F'})_*}"]\\
\pi(\widehat C,F)\arrow[r,"{\pi(\mathrm{id}_{\widehat C},v)_*}"'] \& \pi(\widehat C,F')
\end{tikzcd}}

% foot
% archive sheet number: 4
